\documentclass[11pt]{article}
\usepackage[margin=1in]{geometry}
\usepackage{amsmath,amssymb,amsthm,booktabs,microtype,hyperref}
\hypersetup{hidelinks,pdfauthor={T},pdftitle={Minimax Allocation of Limited Human Review Across Protected Groups}}
\newtheorem{theorem}{Theorem}
\title{Minimax Allocation of Limited Human Review Across Protected Groups}
\author{T}
\date{}
\begin{document}
\maketitle
\begin{abstract}
A team auditing an agent has a fixed number of reviews and may not
know which user group bears the harm. We derive the best worst-group
detection probability for a finite audit with equal review costs.
A continuous calculation is exact because the detection curves are
concave and their fractional allocations admit budget-preserving
randomization. The resulting plan satisfies integer review floors
on every draw. Independent enumeration verifies 128 two-group
instances, including cases where randomization improves on every
fixed allocation. We extend the result to imperfect review and
derive the exact guarantee when miss events can be dependent.
A bound on pairwise joint detections gives a sharp intermediate
curve without specifying the full dependence structure.
\end{abstract}

\section*{Summary and scope}
\paragraph{Question.} How should a service allocate a fixed number of
human reviews when it does not know which group is affected?
\paragraph{Finding.} Concave interpolation gives the exact worst-group
detection optimum. Dependent rounding implements its allocation while
satisfying all integer review floors and the total slot budget on every
draw.
\paragraph{Evidence and scope.} The result is proved and checked
on synthetic queues with independent enumeration. It gives a way to
reserve attention for sparse groups. The guarantee assumes disjoint
groups, equal review costs and stated minimum harm counts.
The imperfect-review extension also needs a lower sensitivity bound
and a declared model for dependence between misses. Pairwise
bounds can tighten the guarantee without assuming independence.

\section{One budget and several places to look}
A service can audit a scheduling or eligibility agent only after
human reviewers inspect the affected cases. A group with few
reports may be small, poorly represented, or simply difficult to
reach. Proportional sampling can miss harm concentrated in such
a group. A separate minimum review count can guarantee that each
group receives attention, but equal counts do not imply equal
chances of finding harm.

Partition a fixed queue into disjoint strata $g=1,\ldots,G$ with
$n_g$ cases. Every review costs one slot and observes a binary harm
label. The unknown affected stratum contains at least $h_g\ge1$
harmed cases if it is affected. Their positions are unknown, and
the reviewer has no reliable case-level signal about them before
inspection. There are $B$ slots and integer floors $a_g$, with
$0\le a_g\le n_g$ and $\sum_g a_g\le B$. Allocate without replacing
a reviewed case. The objective is to maximize the lowest probability
of detecting at least one harm when the affected stratum is chosen
adversarially. Additional affected strata can only help the detection
probability of a fixed randomized allocation.

For a uniform sample of $b$ cases in stratum $g$, the detection
probability at its minimum harm count is
\[
f_g(b)=1-\frac{\binom{n_g-h_g}{b}}{\binom{n_g}{b}},
\]
where the numerator is zero if too few unharmed cases remain.
Let $\bar f_g$ interpolate $f_g$ linearly between adjacent
integers. The result below relies on the equal slot cost and on
that interpolation, rather than on replacing finite sampling
by independent draws with replacement.

\section{An exact continuous allocation}
\begin{theorem}
The best worst-stratum detection probability over all randomized
audit designs is
\[
\max_{\substack{a_g\le x_g\le n_g\\\sum_g x_g\le B}}
       \min_g\bar f_g(x_g).
\]
For any maximizing vector $x$, there is a randomized design
whose integer counts satisfy every floor and the total budget
in every realization. Its expected count in group $g$ is $x_g$,
and its worst-case detection probability in that group is
$\bar f_g(x_g)$.
\end{theorem}

\begin{proof}
First consider the no-harm sequence of observations. For every
fixed choice of the design's random coins, this sequence determines
which cases it would inspect up to budget, including all adaptive
decisions on that sequence. When exactly one stratum is harmed,
the design detects harm if and only if this hypothetical path
intersects its harmed set. Before the first hit all observations
match the hypothetical path. Thus adaptations based solely on
the binary review outcome cannot improve the event of the first
detection beyond the induced random no-hit allocation.

For a given stratum, put its $h_g$ harmed cases uniformly among
its $n_g$ cases for the upper-bound argument. With $b_g$ no-hit
path reviews in that stratum, the detection probability averaged
over placements is $f_g(b_g)$. The worst placement therefore
has detection probability at most $\mathbb E f_g(b_g)$.
Take $x_g=\mathbb E b_g$. The increments of $f_g$ are
nonnegative and nonincreasing. In fact if $q_b=1-f_g(b)>0$,
then the next increment is $q_bh_g/(n_g-b)$.
The ratio of successive positive increments is
\[
\frac{n_g-h_g-b}{n_g-b-1}\le1
\]
whenever both increments are defined and the denominator is
positive. After certainty is reached the increments vanish.
Hence $\bar f_g$ is concave and Jensen's inequality gives
$\mathbb E f_g(b_g)\le\bar f_g(x_g)$. The vector $x$ meets
the continuous constraints, proving the upper bound.

For the converse, write $x_g=m_g+r_g$ with $m_g$ integer
and $0\le r_g<1$. Make consecutive half-open intervals of
length $r_g$ and draw one uniform shift $U$ in $[0,1)$.
Round group $g$ up when its interval contains some point
of $U+\mathbb Z$. An interval of length less than one gets
at most one such point, with probability exactly $r_g$.
The total number of round-ups is at most
$\lceil\sum_g r_g\rceil$. Since the budget is an integer,
the rounded counts sum to at most $B$. The integer floors
and caps are preserved. This is a dependent rounding
construction \cite{gandhi} specialized to one shared budget.

Choose uniformly within each group using its rounded count.
Every fixed harmed set of size $h_g$ then has detection
probability $\mathbb E f_g(b_g)=\bar f_g(x_g)$, because
$b_g$ takes only the two integers adjacent to $x_g$.
Larger harmed sets have at least that probability.
\end{proof}

A target probability $q$ is feasible if the sum of the smallest
$x_g$ with $\bar f_g(x_g)\ge q$ and $x_g\ge a_g$ is at most $B$.
This gives a scalar search. Each inverse curve is piecewise
linear. For rational finite data, finding the breakpoint interval
that contains the target and interpolating gives the exact
optimum.

\section{Where randomization helps}
With one slot and two groups of ten, one harmed case in
the affected group, and no floors, every fixed allocation
leaves one group unchecked. Choose each group with probability
one half. Each has detection probability $1/20$ in its worst
placement. This is a guarantee averaged over the audit draw,
so it does not guarantee that both groups are reviewed on that
draw. Positive floors are the way to impose that separate
requirement when the budget permits it.

For a larger example, take 80 cases with a minimum of eight
harms in the first stratum, 20 cases with a minimum of one in
the second, and ten review slots. The optimal mixed design has
expected counts approximately 3.605 and 6.395. Its worst-stratum
detection probability is $18121/56674$, about 0.31974. The
best fixed allocation gives 0.30. Splitting the slots equally
gives 0.25 and proportional allocation gives 0.10. Here the
smaller stratum receives more reviews because its stipulated
harm density is lower.

The verifier compares the continuous result against independent
two-group enumeration of all feasible allocations and their
line-segment mixtures. In the plane, an optimum of the smaller
coordinate over the convex hull is attained at a vertex or on
an edge, so that oracle is exact. All 128 optima agree.
Randomization strictly improves the fixed allocation in 30
cases, with a largest gap of $2/5$. An independent interval
calculation checks budgets, floors, marginals and achieved
detection on 128 additional three-group cases. These synthetic
counts describe the fixed experiment design.

\section{Imperfect reviews and shared misses}
An inspection can miss a real harm. Suppose each harmed case in
group $g$ has a latent binary mark indicating whether a review
would detect it. The mark law is fixed with the harm placement
before the audit's sampling coins. Clean cases never raise a hit.
The latent marks can represent a shared blind spot in a rubric
or independent reviewer misses. We assume a common lower bound
$s_g\in[0,1]$ on each harmed case's marginal detection probability
conditional on the fixed placement.

\begin{theorem}
If dependence between marks is unrestricted, replace $f_g(b)$
in the minimax allocation theorem by $s_g f_g(b)$. If marks are
independent, replace it by
\[
 f^{\mathrm{ind}}_g(b)
 =\sum_{m=0}^{h_g}\binom{h_g}{m}
   s_g^m(1-s_g)^{h_g-m}
   \left[1-\frac{\binom{n_g-m}{b}}{\binom{n_g}{b}}\right].
\]
In either regime, interpolate the replacement curve and use the
same budget-preserving rounding rule.
\end{theorem}
\begin{proof}
With a uniform sample of size $b$, condition on which cases it
contains. If it contains a harmed case, pick any one of them
using a fixed order on the cases. Its mark is positive with
probability at least $s_g$. Independence of the sampling coins
from the mark law makes the bound valid for this conditional
sample. Thus detection probability is at least $s_g f_g(b)$.
For the matching upper bound, give the $h_g$ harmed cases a
single common Bernoulli mark of mean $s_g$. Use a uniform harm
placement as in the first theorem. The no-hit path analysis
gives exactly $s_g f_g(b)$ after averaging over placements and
the common mark.

For independent marks, the number of detectable harmed cases
stochastically dominates a binomial with parameters $h_g,s_g$.
Conditional on the actual number $m$, uniform review finds a mark
with probability $1-\binom{n_g-m}{b}/\binom{n_g}{b}$. This gives
the displayed lower bound. For the upper bound, take exactly
$h_g$ uniformly placed harms with independent marks of mean $s_g$.
Conditional on their number, detectable cases are uniform
among subsets of that size. The no-hit path then has the
displayed detection curve. Both replacement curves are concave
after interpolation. The first is a nonnegative multiple of
$f_g$, and the second is a mixture of finite-population curves.
Jensen's inequality and the original rounding proof apply.
\end{proof}

The shared-miss model caps guaranteed detection at $s_g$ even
after reviewing the entire group. With independent marks the
cap is $1-(1-s_g)^{h_g}$. The scalar allocation search therefore
starts with the smallest group cap as its upper endpoint.
In a fixture with group sizes 16 and 10, minimum harm counts
five and three, sensitivities $1/2$ and $3/4$, and ten review
slots, the optimal guaranteed detection is $0.4882$ under
arbitrary dependence and $0.6744$ under independence. The
first group receives about $7.328$ expected reviews in the
dependence-robust design and $6.119$ in the independent design.

The verifier checks 1,980 joint-law and slot-count combinations
on one five-case queue. Explicit mark and sample enumeration
checks 100 independent curves. In 96 seeded two-group fixtures,
an allocation oracle confirms both optima, and independence
strictly raises the optimum in 41. These are finite checks of
the stated models. Estimating sensitivity and testing for common
misses require data from the proposed review process.

\subsection{What pairwise detection evidence can buy}
Suppose each harmed case is detectable with probability at least
$s>0$, and each pair is jointly detectable with probability at
most $u$. These are bounds for a fixed harm placement, with the
joint mark law chosen before sampling. Let $M$ be the number of
positive marks among $h$ harmed cases. The bounds imply
$\mathbb E M\ge hs$ and $\mathbb E[M(M-1)]\le h(h-1)u$.

Write $F_b(m)$ for the detection probability under $b$ uniform
reviews when exactly $m$ of $n$ cases have a positive mark. Define
\[
 \mu=hs,\qquad t=\min\{h,1+(h-1)u/s\}.
\]
Let $Z_*$ have mean $t$ and support on the two adjacent integers
around $t$, using a point mass when $t$ is integral. For $m\ge1$ set
\[
 p_m=\frac{\mu}{m}\Pr(Z_*=m),\qquad p_0=1-\sum_{m=1}^h p_m.
\]

\begin{theorem}
The marginal and pairwise bounds are jointly feasible exactly when
$p_0\ge0$. In that case the sharp detection guarantee for every
sample size $b$ is $\sum_m p_m F_b(m)$. The same distribution
$(p_m)$ attains all these minima, and the resulting curve can
replace $f_g$ in the allocation theorem.
\end{theorem}
\begin{proof}
To minimize detection, it suffices to consider $\mathbb E M=\mu$.
For a law with larger mean, independently replace all its marks
with zero often enough to reach $\mu$. This preserves the pairwise
caps and can only lower detection. Define $Z$ by weighting each
positive count by its size, so
$\Pr(Z=m)=m\Pr(M=m)/\mu$. The second-moment bound gives
$\mathbb E Z\le t$. For any function $F$ with $F(0)=0$,
$\mathbb E F(M)=\mu\mathbb E[F(Z)/Z]$.

The sequence $F_b(m)/m$ is decreasing and discretely convex.
To see this for $b>0$, count a hit by giving each selected marked
case a share equal to the reciprocal of the total marks selected.
Conditioning on inclusion of one marked case yields
\[
 \frac{F_b(m)}{m}=\frac bn\,\mathbb E\!\left[\frac{1}{1+K_m}\right],
\]
where $K_m$ counts marked cases among $b-1$ other sampled cases
when $m-1$ of the $n-1$ remaining cases are marked. Couple the
counts for successive $m$ by marking new cases in a fixed order
while retaining the same random sample. The reciprocal decreases.
For two new marks, the expected second difference equals the
expected mixed difference from adding each mark separately, since
interchanging their positions preserves the sample law. It is zero
unless both new marks are sampled. If both are sampled, it is a
nonnegative second difference of $k\mapsto 1/(1+k)$. For $b=0$
the sequence is identically zero.

The linear interpolant of this decreasing convex sequence and
Jensen's inequality give the minimum at the mean $t$, with equality
under $Z_*$. The same argument for $1/m$ shows that $Z_*$ minimizes
$\mu\mathbb E(1/Z)$, the probability of a positive count. Thus
$p_0$ must be nonnegative whenever the bounds are feasible. If it
is nonnegative, sample a count from $(p_m)$ and mark a uniformly
chosen subset of that size. The marginal probabilities are $s$ and
the pairwise probabilities satisfy the cap. This realizes the
claimed curve. With more than $h$ harmed cases the lower bound
still holds by restricting attention to any $h$ of them. The curve
is a mixture of concave finite-population curves, and the uniform
placement and no-hit argument from the first theorem gives the
matching minimax allocation bound.
\end{proof}

With three harmed cases, $s=1/2$ and $u=1/4$ give $t=2$. The
least favorable law has two marks with probability $3/4$ and
zero marks otherwise. A full review therefore guarantees detection
with probability $3/4$. The corresponding guarantees are $1/2$
under unrestricted dependence and $7/8$ under independence.
The pair cap here is the same as the joint probability under
independence, yet even exact pairwise probabilities leave room
for a shared blind spot across all three cases.

The exact verifier enumerated the vertices of the count moment
polytope for 1,296 panels. It found 274 infeasible panels and
matched 6,682 sample-size minima on the feasible ones. An audit
can use a joint-detection bound derived from adjudicated paired
cases if the design justifies transferring that bound to the queue.
The formula supplies the detection guarantee conditional on that
bound. It does not estimate it from unaudited cases.

\newpage
\section{When the model needs to change}
Unequal review costs can invalidate the interpolation. For
example, let there be one case per group, review costs two and
one, and budget one. A fractional calculation with expected
counts $(1/3,1/3)$ satisfies the expected budget but no feasible
realization can inspect the first group. Its worst-group
detection probability is zero. In that setting a mixed strategy
must be supported on actually feasible cost allocations.
Overlapping group membership also changes the feasible set
because one review can count toward several floors.

A prospective audit should preregister group definitions,
actual slot costs, and the harm counts for which it reports
power. The method can balance scarce review effort for an
agent's decisions, while a reliable predictor, variable review
duration, or false positive labels would change the optimization.
When review sensitivity depends on assignments or on prior
reviews, evaluate that process before adopting either mark law.
Report the realized allocation, the audit design's detection
probability, and any failed coverage floors as separate facts.

\paragraph{Review status.} Independent manuscript review is pending.

\begin{thebibliography}{9}
\bibitem{gandhi}
Rajiv Gandhi, Samir Khuller, Srinivasan Parthasarathy, and
Aravind Srinivasan. Dependent Rounding and Its Applications to
Approximation Algorithms. Journal of the ACM 53(3), 2006.
\url{https://www.cs.umd.edu/~srin/PDF/2005/depround-dec05.pdf}.
\end{thebibliography}
\end{document}
